\chapter{Internal SOPs}
\label{ch:builder-internal-sops}

\section{What will you build in this chapter?}
\label{sec:builder-internal-sops-intro}

An internal SOP is a workflow of consecutive steps that starts on an event you choose, such as a record being created or a form being submitted. Its steps may request an approval, create a task or send a notification, so getting the work done does not depend on someone remembering the next step. This chapter explains how to build an SOP, test it, publish it and follow up its runs.

\begin{rcnote}[Example: Horizon]
When a new visit request is created, an internal SOP requests Sarah Al-Khatib's approval, and once she approves it creates the visit preparation task for Omar Haddad.
\end{rcnote}

\section{Pages in this chapter}
\label{sec:builder-internal-sops-pages}

You start from the “Internal SOPs” page in Builder's “Core” section.

\begin{table}[H]
\centering
\caption{Running work with internal SOPs}
\label{tab:builder-pages-internal-sops}
\rcZebra
\begin{tabularx}{\linewidth}{R{4.2cm} Y}
\rcHeaderRow
\thh{Page and section} & \thh{What this page offers}\\
Internal SOPs — Core & Builds workflows of steps, such as approvals and task creation, that start on an event you choose.\\
\end{tabularx}
\end{table}
